\section{Warm starts}\label{supp:warm}

\paragraph{An iteration bound}
A prediction can serve as the initial iterate of the conjugate gradient method
\cite{HestenesStiefel1952} applied to $K\Ustar=F$. Let $u^k$ be the iterate after $k$
iterations and $\gap_k=\gap(u^k)$, so that $\gap_0$ is the energy gap of the initial
iterate. For $\kappa>1$ the classical bound \cite[Section~6.11.3]{Saad2003} is
\begin{equation}
\frac{\gap_k}{\gap_0} \le 4q^{2k},
\qquad
q = \frac{\kappa^{1/2}-1}{\kappa^{1/2}+1}.
\label{eq:cheb}
\end{equation}

\begin{corollary}[Iteration bound of a warm start]\label{cor:warm}
Let $\bar g\in(0,1)$ be a target relative energy gap and $g_0\in(\bar g,1]$ the relative
energy gap of the initial iterate. Under \eqref{eq:cheb}, $k\ge\ln(4g_0/\bar g)/(2\ln
q^{-1})$ iterations suffice to reach $\bar g$, and $g_0=1$ for the zero field. The warm
start therefore lowers this iteration bound by $\ln(1/g_0)/(2\ln q^{-1})$ iterations, a
fraction
\begin{equation}
\frac{\ln(1/g_0)}{\ln(4/\bar g)}
\label{eq:warm}
\end{equation}
of the bound for the zero field; the fraction does not depend on $\kappa$.
\end{corollary}

\begin{proof}
By \eqref{eq:cheb}, $g(u^k)=g_0\gap_k/\gap_0\le4g_0q^{2k}$, which is at most $\bar g$ once
$k$ reaches the stated value; subtract the two iteration counts.
\end{proof}

\Cref{cor:warm} states in the energy norm an estimate that Sun et~al.~\cite{Sun2026FEVessel}
give in the relative residual, after a bound of Wang et~al.~\cite{Wang2026PretrainFEM} for
fixed-point iterations. It compares bounds, not iteration counts: how many iterations a warm
start actually saves depends on the spectral content of its error and can even be negative.
What it shows is that a prediction of given accuracy lowers the iteration bound by a fixed
number of iterations, while the bound grows with $\ln(1/\bar g)$, so that its share shrinks
as the tolerance tightens. Sun et~al.~\cite{Sun2026FEVessel} draw the same conclusion for
conjugate gradients in the relative residual, where the constant of the bound and the
conversion between residual and error add a bounded offset to the saving, and Wang
et~al.~\cite{Wang2026PretrainFEM} from an asymptotic bound of the same form for fixed-point
iterations; \cref{cor:warm} states it in the energy norm, without that offset, and
\cref{rem:residual} bounds the initial residual by the energy gap.

\begin{remark}[Residual-based stopping]\label{rem:residual}
The initial residual of a prediction is $r_0=-Ke_0$, so
$\norm{r_0}_2^2=e_0^{\top}K^2e_0\le\lambda_{\max}\norm{e_0}_K^2$, and with
$\norm{\Ustar}_K^2=F^{\top}K^{-1}F\le\norm{F}_2^2/\lambda_{\min}$,
\begin{equation}
\frac{\norm{r_0}_2}{\norm{F}_2} \le (\kappa\,g_0)^{1/2},
\label{eq:residual}
\end{equation}
while the zero field starts from $\norm{r_0}_2/\norm{F}_2=1$. Both inequalities are
attained, the first when $e_0$ is an eigenvector of $\lambda_{\max}$ and the second when $F$
is one of $\lambda_{\min}$, so the bound is sharp. A prediction that is accurate in energy
can therefore start with a larger residual than the zero field, by a factor of up to
$(\kappa g_0)^{1/2}$, because $K$ amplifies the stiff-mode components of its error, and a
residual-based stopping test then credits a warm start with less than its energy accuracy.
Wang et~al.~\cite{Wang2026PretrainFEM} observed the same for Newton's method: the initial
residual of their predictions could exceed that of the zero initial guess although their
errors were much smaller.
\end{remark}

\paragraph{The tests}
For the prediction tested below, with an initial relative energy gap
$g_0=\numKfiveStartGap$, the fraction \eqref{eq:warm} is \numWarmFractionSeven{} for $\bar
g=10^{-7}$, \numWarmFractionTwelve{} for $\bar g=10^{-12}$ and \numWarmFractionSeventeen{}
for $\bar g=10^{-17}$. A residual test $\norm{r_k}_2\le\tau\norm{F}_2$ with $\tau=10^{-6}$
corresponds to targets in this range: it guarantees $g\le\kappa\tau^2$, which is up to about
\numKappaTauTwo{} for the two-dimensional condition numbers of \cref{sec:methods:data}, and
an iterate exactly at the threshold has $g\ge\tau^2/\kappa$, which is down to about
\numTauSqKappaTwo{}.

The July run tested a warm start under criterion K5. On \numKfiveEval{} validation instances
and four load cases each, unpreconditioned conjugate gradients needed on average
\numKfiveZero{} iterations from the zero field to reach a relative residual of
\numKfiveTol{}, \numKfiveFree{} from the prediction of a label-free transformer of the
earlier code, trained for 100 epochs for this test, and \numKfiveNaive{} from the scaled
polynomial: savings of \numKfiveSaving{} and \numKfiveNaiveSaving{} against the 20\%
required, so that K5 was triggered and the claim withdrawn. The prediction had a relative
energy gap of \numKfiveStartGap{}, which five and twenty iterations from it lowered to
\numKfiveFiveGap{} and \numKfiveTwentyGap{}, but it did not shorten the run to the
tolerance. This is consistent with \cref{cor:warm,rem:residual}: at a tight tolerance a warm
start of this accuracy lowers the iteration bound by only a small fraction, and a
residual-based stop can see a prediction accurate in energy as further from convergence than
the zero field. Neither the initial residuals nor energy-based stopping were measured, so
the two effects cannot be separated here. The test used one network of the earlier code,
trained for half the usual number of epochs, and one solver without preconditioning. The
predictions timed in \cref{tab:cost}, started instead of the zero field, needed between
\numCostWarmMore{} more and \numCostWarmFewer{} fewer iterations of conjugate gradients to a
relative residual of \numCostCgTol{} (medians per set), again consistent with
\cref{cor:warm,rem:residual}; their initial residuals were not recorded either.

\paragraph{Reported speed-ups}
The pretrain finite element method needed up to about nine times fewer conjugate-gradient
iterations than a zero start at a tolerance of $10^{-3}$ \cite{Wang2026PretrainFEM}. In its
pressure-vessel follow-up the saving of conjugate gradients fell from about three times at a
relative residual of $10^{-3}$ to almost none at $10^{-4}$ \cite{Sun2026FEVessel}, and
neural-operator warm starts have also been reported to accelerate iterative solvers
\cite{Eshaghi2026NOWS}. Wang et~al.~\cite{Wang2026PretrainFEM} and Sun
et~al.~\cite{Sun2026FEVessel} explain why the benefit shrinks as the tolerance tightens: a
warm start lowers an iteration bound by a fixed number of iterations, as Wang et~al.\ show
for fixed-point iterations and Sun et~al., in the relative residual, also for conjugate
gradients. The bound is not the count: with conjugate gradients, the warm start of Sun
et~al.\ saved 580 iterations at $10^{-3}$ but 69 at $10^{-4}$ (their Table 3).
\Cref{cor:warm} states the bound for conjugate gradients in the energy norm: a warm start
lowers the iteration bound by a number of iterations set by its own accuracy, so the share
of the run it can save shrinks as the tolerance tightens, and a residual-based stop can hide
part of it (\cref{rem:residual}). Reported warm-start speed-ups are therefore comparable
only when the tolerance, the stopping test and the preconditioner are stated with them: Wang
et~al.~\cite{Wang2026PretrainFEM} report, for their linear problems, iteration counts to a
relative residual of $10^{-3}$, with unpreconditioned conjugate gradients in two dimensions
and with unpreconditioned and preconditioned conjugate gradients in three, and Sun
et~al.~\cite{Sun2026FEVessel} report iteration counts and wall-clock speed-ups at three
relative-residual tolerances, the speed-ups computed from times that include the network's
forward pass and, for algebraic multigrid, its set-up.
